\documentclass[letterpaper]{article}

\usepackage[letterpaper,margin=1in]{geometry}

% Compile from the repository root so that this path resolves to
% the pinned WSP submodule.
\input{wsp/testing/test-case-library.tex}

\begin{document}

\def\TCID{%
  TC-0012
}

\def\TCTitle{%
  Toolbar buttons
}

\def\TCPurpose{%
  Verify the buttons of a Toolbar: styles, button menus, the group rule,
  errors, SaveToolbar and RestoreToolbar, and the events.
}

\def\TCPriority{%
  \TCPriorityHigh
}

\def\TCRequirementRef{%
  REQ-0012, REQ-0020, REQ-0022
}

\def\TCDesignRef{%
  \texttt{ADR-0006} and \texttt{ADR-0010} in WFC
}

\def\TCFeatureRef{%
  \texttt{tests/corpus/190-comctl-toolbar.frm} and
  \texttt{tests/corpus/190-comctl-toolbar.vbp}; the form code and the controls
  of \texttt{include/wfc/mscomctl/mscomctl.inc} and the runner
  \texttt{tests/tc-0012-comctl-toolbar.ps1}
}

\def\TCTestTechnique{%
  Scenario-based testing of one form: each Immediate-window line is compared
  with the recorded outcome (ISO/IEC/IEEE 29119-4 experience-based and
  specification-based techniques, with expected output comparison).
}

\def\TCPreconditions{%
  A \texttt{wfc} of the same version as the library (WFC, the interpreter and
  compiler these controls belong to, runs the case: the library holds the
  program, WFC the runner). The files of the case are in
  \texttt{tests/corpus}: \texttt{190-comctl-toolbar.vbp},
  \texttt{190-comctl-toolbar.frm}, \texttt{190-comctl-toolbar.expected},
  \texttt{190-comctl-toolbar.debug}, \texttt{190-comctl-toolbar.ui}.
}

\def\TCEnvironment{%
  Windows with \texttt{wfc.exe} of the x64 Debug build; the form runs in WFC's
  headless host (no desktop is needed). WFC registers the case with CTest and
  runs \texttt{wfc --interpret 190-comctl-toolbar.vbp}.
}

\def\TCAssumptions{%
  The scripted input \texttt{190-comctl-toolbar.ui} is the only input; the
  case does not read the registry or other files except the picture named in
  the form.
}

\def\TCInputData{%
  A form with a Toolbar; the buttons \texttt{new}, \texttt{sep}, \texttt{bold}
  (check), \texttt{left} and \texttt{center} (group) and \texttt{drop}
  (drop-down with two menus); the script clicks each button and opens the menu
  \texttt{m2} and the customization.
}

\def\TCInitialState{%
  A fresh \texttt{wfc} process; the program creates its own controls from the
  form file and starts with no state from a previous run.
}

\def\TCProcedure{%
  \begin{enumerate}
    \item Add the six styles of button and two button menus; print the counts,
      the menu text, the style and the type names.
    \item Print the separator's width, a button's position and the top of the
      first button.
    \item Press \texttt{left} then \texttt{center} and print both values.
    \item Set style 9 and set \texttt{Left} and print the errors.
    \item Save the toolbar, hide \texttt{new}, restore it and print
      \texttt{Visible}.
    \item Run the script.
  \end{enumerate}
}

\def\TCExpectedResult{%
  \begin{enumerate}
    \item Six buttons, two menus, \texttt{Recent 2}, style 5, type names
      \texttt{Button} and \texttt{ButtonMenu}; the separator is 8 pixels wide
      and the buttons follow it from left to right.
    \item Pressing \texttt{center} releases \texttt{left}.
    \item A style of 9 is 380 and writing \texttt{Left} is 383 (read only).
    \item Restoring brings \texttt{new} back (\texttt{True}).
    \item Events: \texttt{click new 0}, \texttt{click bold 1}, \texttt{click
      bold 0}, \texttt{click left 1}, \texttt{dropdown drop}, \texttt{menu
      m2Recent 2}, \texttt{changed}.
    \item \texttt{wfc} exits with code 0, writes the exact text of
      \texttt{190-comctl-toolbar.expected} to standard output, and writes the
      exact Immediate-window text of \texttt{190-comctl-toolbar.debug} to
      standard error. Any difference fails the case.
  \end{enumerate}
}

\def\TCPostConditions{%
  The form is closed by the last \texttt{close Form1} line of the script;
  nothing is left running and no file is written.
}

\TCTemplate

\end{document}
